\section{Introduction}

Rare diseases affect approximately 350 million individuals globally across nearly 7,000 distinct conditions, yet each condition represents less than one percent of the overall disease burden \cite{finetti2025data}. This epidemiological paradox creates substantial barriers to clinical research and diagnostic development. Patients experience an average diagnostic delay of six years from symptom onset to accurate diagnosis, during which time they navigate fragmented healthcare systems with limited access to specialized expertise. The combination of small patient cohorts, geographic dispersion, and phenotypic heterogeneity results in datasets ranging from dozens to a few thousand cases---far below the sample sizes required for robust machine learning model training \cite{finetti2025data,mendes2025synthetic}.

Traditional approaches to data augmentation, including geometric transformations and photometric adjustments, have dominated biomedical applications due to their interpretability and computational efficiency. A recent scoping review of 118 studies published between 2018 and 2025 found that classical augmentation techniques accounted for 49.3\% of data expansion strategies in rare disease research \cite{finetti2025data}. However, these methods are fundamentally constrained by the variability present in the original dataset, limiting their capacity to generate novel yet biologically plausible samples. As artificial intelligence applications in healthcare have expanded, the need for more sophisticated data generation approaches has become increasingly apparent.

Deep generative models have emerged as a promising solution to overcome the limitations of classical augmentation. Generative Adversarial Networks (GANs), which employ adversarial training between generator and discriminator networks, have demonstrated particular success in medical imaging applications. Frid-Adar et al. achieved a 7.1 percentage point improvement in liver lesion classification sensitivity through GAN-based synthetic CT image generation, with final performance reaching 85.7\% sensitivity and 92.4\% specificity \cite{fridadar2018gan}. The application of GANs has extended beyond imaging to electronic health records (EHRs), where medGAN successfully synthesized multi-label discrete patient records with 97.77\% accuracy while preserving privacy through fully synthetic data generation \cite{choi2017generating}.

Variational Autoencoders (VAEs) offer an alternative approach through probabilistic modeling of latent distributions, demonstrating particular utility in scenarios with limited computational resources. While VAEs may produce slightly less sharp images compared to GANs, their lower computational requirements and reduced risk of mode collapse make them valuable for resource-constrained settings \cite{mendes2025synthetic}. More recently, diffusion-based models have introduced novel capabilities for sequential data generation. The MedDiffusion framework, applying denoising diffusion probabilistic models to EHR data for the first time, achieved 77.88\% PR-AUC on kidney disease risk prediction by generating synthetic patient visits conditioned on medical history \cite{zhong2024meddiffusion}.

Despite these advances, substantial challenges remain. Finetti et al. identified that over 40\% of published studies lack sufficient quantitative detail regarding dataset sizes before and after augmentation, hindering reproducibility and comparative analysis \cite{finetti2025data}. Privacy concerns persist even with synthetic data generation, as membership inference and model inversion attacks remain potential vulnerabilities \cite{mendes2025synthetic}. Validation practices vary widely, with most studies evaluating synthetic data quality solely through downstream model performance rather than independent biological plausibility assessment. The computational demands of deep generative models, requiring GPU clusters and extensive hyperparameter tuning, create additional barriers to widespread adoption.

This review synthesizes current evidence on generative model applications for rare disease biomedical data synthesis. We examine the technical foundations of GANs, VAEs, and diffusion models in Section~\ref{sec:methods}, analyze strategies for mitigating data scarcity and class imbalance in Section~\ref{sec:scarcity}, review modality-specific applications across imaging, EHRs, and omics data in Section~\ref{sec:applications}, discuss current limitations in Section~\ref{sec:limitations}, and conclude with future directions for clinical integration in Section~\ref{sec:conclusion}. Our focus remains on how these generative approaches specifically address the unique constraints imposed by rare disease research, where traditional data collection strategies fail to accumulate sufficient samples for modern machine learning pipelines.

